\chapter{Аудит кандидатов абсолютного масштаба такта рождения}
\label{ch:v10-nonequilibrium-bath-birth-tick-scale-audit}

Абсолютный кандидат на $\tau_{\rm birth}$ должен удовлетворять шести
условиям: иметь размерность времени, присутствовать во внутреннем носителе,
выбираться существующим родителем, независимо разрушать масштабную орбиту,
не использовать наблюдаемую цель и иметь однозначный типизированный морфизм
в такт рождения.

Проверены двенадцать классов: $\hbar/E_C$, световое время ребра,
UV-фаза, K43-обрезание, время корреляции ванны, хопфовская проводимость,
dimensional transmutation, планковское и кривизненное времена, наблюдаемый
темп роста, безразмерное действие вакуума и внешний атомный стандарт.
Матрица размера $12\times6$ имеет ранг $6$, но полных проходов нет.

Особенно существенно, что пять лучших внутренних записей не независимы:
\begin{equation}
 \frac{\hbar}{E_C}
 =\frac{\ell_{\rm edge}}c
 =\frac{2\sqrt3}{\omega_{\rm UV}}
 =\frac{42}{c\Lambda_{43}}
 =\frac1{22\kappa}
 =\tau_{\rm birth}.
\end{equation}
Каждая получает $5/6$, но всем недостаёт одного и того же критерия ---
независимого разрушения масштабной орбиты. Переписывание одной формулы через
другую не создаёт нового якоря.

Время корреляции также не закрывает задачу. При
$\omega_{\rm UV}\tau_{\rm birth}=2\sqrt3$ экспоненциальный и гауссов
профили дают соответственно
\begin{equation}
 \frac{\tau_E}{\tau_{\rm birth}}=\frac1{2\sqrt3},
 \qquad
 \frac{\tau_G}{\tau_{\rm birth}}=\frac{\sqrt\pi}{4\sqrt3},
\end{equation}
то есть различаются уже до появления абсолютной единицы.

Планковское, кривизненное, наблюдаемое и атомное времена действительно могли
бы повысить ранг карты с $3$ до $4$, но первые три требуют невыведенных или
круговых физических величин, а атомный период является явным внешним
эталоном. Безразмерное $S_{\rm vac}$ вообще не имеет временной размерности.

\begin{theorem}[Исчерпание текущих кандидатов такта]
Двенадцать классов покрывают все шесть критериев с полным критериальным
рангом, однако ни один внутренний кандидат не разрушает общую орбиту
$(\tau,E,\ell)\mapsto(s\tau,E/s,s\ell)$. Число полных проходов равно нулю.
\end{theorem}

\begin{nogocriterion}[Эквивалентные часы нельзя считать независимыми]
Формулы, связанные уже установленными родительскими равенствами, являются
разными координатными представлениями одного тика. Их совместное перечисление
не повышает размерный ранг и не может служить доказательством абсолютного
времени.
\end{nogocriterion}

\begin{protocol}\begin{enumerate}
\item проверенные кандидаты: \textbf{$12/12$};
\item критериальный ранг: \textbf{$6/6$};
\item полные проходы: \textbf{$0/12$};
\item внутренние разрушители масштабной орбиты: \textbf{$0$};
\item происхождение абсолютного такта: \textbf{$0/1$};
\item происхождение физического часового носителя: \textbf{$0/1$};
\item следующий гейт: \textbf{dimensional transmutation такта рождения}.
\end{enumerate}\end{protocol}

Машинный сертификат сохранён в
\path{s2t_v10_cell_birth_four_volume_nonequilibrium_bath_birth_tick_absolute_scale_candidate_audit_gate_results.json}.